\documentclass[article,twoside,11pt]{report}



%Images
\usepackage[pdftex]{graphicx}

\usepackage[usenames,dvipsnames]{xcolor}

%In-document links
\usepackage[linktoc=all, colorlinks=true, linkcolor=black, urlcolor=blue, citecolor=black]{hyperref}

%headers and footers
\usepackage{fancyhdr}

\usepackage[margin=1.4in]{geometry}





%text representing commands or in-program text
\newcommand{\command}[1]{\texttt{#1}}
%Place an icon inline with the text
\newcommand{\icon}[1]{\includegraphics[height=1em]{icons/#1.png}}
%Place and icon and command text together to represent a button
\newcommand{\button}[2]{\ \command{\icon{#1} #2}}
%Keyboard shortcut combining a modifier key and another key
\newcommand{\shortcut}[2]{\command{#1 + #2}}
%Menu chain
\newcommand{\menu}[0]{$\rightarrow$}

%filesystem path
\newcommand{\file}[1]{\command{#1}}

%code snippets
\newcommand{\code}[1]{\command{#1}}
\newcommand{\class}[1]{{\color{violet} \code{#1}}}

\definecolor{commentcolour}{rgb}{0.1, 0.5, 0.2}
\newcommand{\comment}[1]{{\color{commentcolour} \code{#1} }}

\definecolor{codeblockbackground}{rgb}{0.9, 0.9, 0.9}
\newcommand{\codeblock}[1]{{%
	\vspace{1em}%
	\centering%
	\fcolorbox{white}{codeblockbackground}{%
		\parbox[l]{\textwidth}{%
			
			\code{#1}%
		}%
	}%
	\vspace{1em}%
}}

\newcommand{\tab}{\hspace*{2em}}



%element
\newcommand{\element}[1]{$#1$}

%emphasis
\newcommand{\emphasis}[1]{{\bf #1}}

%placing a screenshot as a figure
\newcommand{\screenshot}[2]{%
\begin{figure}[h!]
\centering\includegraphics[width=0.85\textwidth]{figures/#1.png}
\caption{#2}
\end{figure}
}


%Table of Contents macros
\newcommand{\tocchapter}[1]{\cleardoublepage\chapter*{#1}\addcontentsline{toc}{chapter}{#1}}
\newcommand{\tocsection}[1]{\section*{#1}\addcontentsline{toc}{section}{#1}}
\newcommand{\tocsubsection}[1]{\subsection*{#1}\addcontentsline{toc}{subsection}{#1}}
\newcommand{\tocsubsubsection}[1]{\subsubsection*{#1}\addcontentsline{toc}{subsubsection}{#1}}

\title{Peakaboo 4}
\date{March 6, 2012}
\author{Nathaniel Sherry, Marina Suominen Fuller}
        

\begin{document}

\renewcommand{\headrulewidth}{0pt}
\lhead[]{}
\rhead[]{}
%\rhead[\includegraphics[height=1em]{title/peakaboo.pdf}]{\includegraphics[height=1em]{title/peakaboo.pdf}}
\fancyhead[C]{\includegraphics[height=1em]{title/peakaboo.pdf}\ \includegraphics[height=1em]{title/usersguide.pdf}}


\pagestyle{empty}



\begin{titlepage}
	\vspace*{\fill}
	\centering
	\includegraphics[height=4em]{title/peakaboo.pdf}\\
	\vspace*{1em}
	\includegraphics[height=1.5em]{title/usersguide.pdf}\\
	\vspace*{\fill}
	\vspace*{\fill}
\end{titlepage}

\cleardoublepage

\pagestyle{empty}
\setcounter{tocdepth}{1}
\tableofcontents
\addtocontents{toc}{\protect\thispagestyle{empty}}
\cleardoublepage
\setcounter{page}{1}
\pagestyle{fancyplain}


\tocchapter{Introduction}

Peakaboo is a part of the Science Studio package that allows users to identify the 
spectral origins of an XRF spectrum using a routine that fits all components of 
the K, L, or M spectrum including escape peaks and pileup peaks, and then plots 
their spatial intensity distributions as maps.

\tocsection{Accepted File Formats}

Peakaboo supports several file formats for XRF data, including: 

\begin{enumerate}

\item \command{CDFML}, an xml-based file format defined by the Canadian Light 
Source using NASA's CDF/CDFML container.

\item \command{Peakaboo Plaintext Format}, a simple XRF format comprised of 
rows of space-separated numbers. Each row represents a single spectrum.

\item \command{CLS Data Acquisition Format} used by the Canadian Light Source 
for XRF collection on the VESPERS beamline.

\end{enumerate}

For a complete listing of supported file formats, see ``File Format Descriptions''
in the appendix.

Peakaboo also allows users who are familiar with the Java programming language to 
add support for new file formats. See the section on ``Extending Peakaboo'' for more
information.

\tocsection{Getting Help}

If you have any questions or concerns not addressed in this guide, please feel free
to contact the developers at \command{nsherry4@uwo.ca}


\tocchapter{Getting Started}

\tocsection{Opening Data}

Select \button{document-open}{Open} from the toolbar or 
\command{File \menu \button{document-open}{Open}} from the menu to locate your 
XRF data. If all scan point data is contained in a single file, simply select
that file and open it. If there are many individual scan point files, 
(eg one scan per file), then locate them in their directory and use 
\shortcut{Control}{A} to select all of them.

\screenshot{open-files}{Peakaboo's file selection dialog}

\tocsection{File Format Conflicts}

While rare, it is possible that more than one of the modules that Peakaboo uses to read XRF data
will believe that it can open a particular data set. This is usually the result of many
different file formats having some of the same properties, such as similar file 
extensions (eg \file{.dat}, \file{.xml}, \ldots ). In the event that Peakaboo cannot automatically 
determine which XRF reader module to use, it will ask you to indicate which format your
data is in. Select the entry the description of which best matches your data.

\screenshot{select-datasource}{Peakaboo's file format conflict dialog}

\tocsection{Examining the Data}

The \icon{zoom-in} \icon{zoom-out} zoom slider bar on the bottom right of the window 
allows you to expand the x-axis (energy scale). You can use the scroll bar above
this area, click and drag on the spectrum itself, to pan back and forth to view 
different energy regions.

\screenshot{zoom-scroll}{Zoom controls and horizontal scrollbar}


\tocsection{Individual Scans}

The \command{Scan \#} field on the left side, under the spectrum view, allows you 
to view individual spectra in a map. Click on the arrows to move in the desired 
direction or enter a value and press \command{Enter}. This option will be disabled 
if you are viewing the mean average or strongest per channel composites rather than 
individual scans.

\screenshot{scan-number}{Scan number selector}


\tocsection{Energy Levels}

The \command{Max Energy} setting on the upper right provides energy calibration. 
If your dataset contains this value, the maximum energy value will be set 
automatically. If not, you will have to enter it manually. This value can be 
tweaked or adjusted by selecting a major element known or suspected to be present 
in the sample, and adjusting the energy range until the associated peak fits properly. 
First add your axes to the spectrum by selecting \command{View \menu\ Axes} from 
the menu. The Y Axis -- \command{Relative Intensity} and X Axis -- \command{Energy (keV)} should now be 
labelled for you.

\screenshot{max-energy}{The maximum energy selector}

On the left side of the screen, under the \command{Peak Fitting} tab, click on the 
drop-down arrow beside the \button{edit-add}{Add Fittings} button and select 
\command{Elemental Lookup}, and a list of elements will become visible, 
listed in increasing atomic number. Find the element you wish to use as a reference, 
and click on the checkbox on the box next to it. A set of peaks will appear in shaded 
red within the spectrum. Adjust the Max Energy (keV) setting to move the red shaded 
peaks to fit these intense peaks in the spectrum. The Max Energy setting should now
be calibrated to your data.

You should double-check this Max Energy calibration by adding another element. Although
\element{Ar} may not be present in your sample, it typically appears in XRF spectra.
Select \element{Ar} from the peak list and the \element{Ar} K-$\alpha$ peak. It should appear 
at 2.95 keV with the accompanying K-$\beta$ peak. If you start adding element lines and the fits
appear off, check your calibration again.

\tocsection{Customizing the View}

The \command{View} menu contains many settings to help you to visualize the 
spectrum: \command{Logarithmic Scale} (usually best), \command{Axes} (shows energy 
and intensity scales), \command{Mean Average of Spectra} is the average of all scans 
(averages each energy channel across all spectra within a map to give you a much 
improved composite spectrum).

You can also add curve fitting markings to your spectrum: 
\command{View \menu\ Curve Fit \menu\ Element Names} to label each curve with the 
name of the element it represents,
\command{View \menu\ Curve Fit \menu\ Markings} showing transition markings (K-, L-lines), 
and \command{View \menu\ Curve Fit \menu\ Heights} for fitting heights (max counts).

\screenshot{view-example}{Peakaboo with mean averaged data, axes and fitting labels}

Your spectrum, as it appears in the window, can also be saved by selecting 
\button{device-camera}{Export Plot as Image} from the toolbar, 
\command{File \menu\ Export Plot as Image}, or \shortcut{Control}{P}.  
Image format options include \command{PNG}, \command{SVG} or \command{PDF}. You can 
also save the fittings results as a text file by selecting 
\command{File \menu \button{document-export}{Export Fittings as Text}}.


\tocsection{Saving and Loading Sessions}

You can save your session information, so that next time you want to look at your XRF 
data, you can just load this information (session) and not have to start from the 
beginning. These options are found under \command{File \menu \button{document-save}{Save Session}} 
and \command{File \menu\ Load Session}.




\tocchapter{Peak Fitting}


\tocsection{Escape Peaks}

Before doing any peak fitting, you should consider the type of
detector you used to collect the XRF spectra. Features may be present in your
spectra that are due to incoming X-rays interacting with the detector material. For
example, for a \element{Si} detector, there is a probability that some of the incoming X-rays
will interact with the \element{Si} and kick out \element{Si} K-shell electrons, thereby reducing the
incoming X-ray's measured energy by 1.74 keV. This escape peak is present for the
major elements in the sample.

To include escape peaks in fitting the peaks in your spectra, select 
\command{View \menu\ Escape Peaks} and choose either \command{Silicon} or 
\command{Germanium}, depending on the type of detector you used. 
If you do not need to fit escape peaks, choose \command{None}.

\tocsection{Selecting Fittings}

There are three different tools which can be used to fit your data. To access these 
options, click on the drop-down arrow next to the \button{edit-add}{Add Fittings} button at 
the top of the \command{Peak Fitting} tab.

\screenshot{fitting-lines}{Detail of all fitting lines for Iron K}

\tocsubsection{Guided Fitting}

This is the default option for \button{edit-add}{Add Fitting}. Click on the
peak you wish to fit and the program will fit the peak with a possible option. Other 
possible element lines will also be displayed in the drop down list. To add another 
fitting, click on the \button{edit-add}{Add}  button below the current entry and 
click on another peak. To edit a previous fitting, click \button{edit-edit}{Edit}. 
Once you are happy with the fittings click \button{choose-ok}{OK} to add the 
fittings to your list.

\screenshot{guided-fitting}{The guided gitting controls}


\tocsubsection{Element Lookup}

Best results are usually obtained by adding elements with the lowest atomic numbers 
first and working your way up. It is recommended that you fit the K lines first. 
L lines for elements for tungsten and higher should be fit separately from the K 
lines for lower atomic number elements, as the K lines will interfere with the L 
line fittings. As you select elements, their fittings will appear in red on the 
spectrum. Once you are happy with the fittings, click \button{choose-ok}{OK}, and the
list of newly-fitted elements will be added to your fitted element list, and will 
appear in black on the spectrum.

\screenshot{lookup-fitting}{The elemental fitting controls}

\tocsubsection{Summation Peaks}

If you are having trouble finding element lines to fit some peaks in your spectrum,
you may have what are called `pile-up peaks' (referred to as `summation' peaks in
Peakaboo). This is a detector phenomenon. If your sample has a lot of one element
present, then some of the emitted x-rays due to the presence of this element will
impinge on the detector virtually simultaneously and the pulse created and
measured would be the sum of the two x-ray energies. Please note that for element 
lines to appear in the drop down lists, they must have already been included in the 
spectral fit.

\screenshot{fitting-summation}{The summation fitting controls}

\tocsection{Working with Fittings}

Each fitting has a checkbox under the \command{Fit} column in the \command{Peak Fitting}
panel. This checkbox controls whether this fitting is enabled or not. Toggling the 
checkbox will allow you to temporarily remove a specific fitting from the fitting process. 
Rechecking the box will restore the fitting as it was.

When using background filters, it is usually best to leave 10\%-15\% of the background in place
(default is 10\%) in order to prevent the background filter from removing all signal under the
edges of the peak. When too much signal around the edges of a peak is removed, it can prevent
peaks from being fitted to the proper intensity, or being fitted at all.

Ordering of fittings is also important, as fittings listed first are the first 
to be fit to the data. While there are usually multiple peaks per fitting which all constrain 
the height of a fitting, in the case of peak overlap, a fitting will occasionally 
be calculated to be larger than it ought to be. In this case, other fittings may be 
underrepresented because there is no more signal left for those fittings to be fit to.
Changing the order of fittings can reveal and correct these problems. Select a fitting 
(\shortcut{Control}{P} to select several) that you wish to change the ordering of, and
click \button{go-up}{Move Up} or \button{go-down}{Move Down} to adjust the ordering
of their fitting.

Adjusting the \command{Max Energy} setting by $\pm$ 0.05 keV can also sometimes improve fitting
if the energy value provided with the data is not exactly accurate. See the section on Energy 
Levels under ``Getting Started'' for more details.

\tocsection{Details of Peak Fitting}

A fitting is a model composed of several ``transitions'' from one of the K, L, or M groups of 
electron transitions. The positions and relative intensities for each transition series
were taken from several tabulated sources. Fittings take account of details such as 
separations and relative intensities of transitions such as the K$\alpha$ 1 and 2,  
K$\beta$ 1, 2, and 3, L$\alpha$, L$\beta$ 1 and 2, and the L$\gamma$ 1, 2, and 3 transitions.

For most of the K series elements from \element{Ca} to \element{Mo}, the relative 
intensities of alpha and beta composite transitions were checked and adjusted for our 
lookup table using metals or compounds. As well, a number of metals with strong L 
series spectra were analysed and the relative intensities adjusted accordingly; 
most (but not all) ratios were found to agree well with the tabulated sources. The 
intensity rations for the K transitions have been optimized for $\sim$20KV irradiation; relative 
intensities for most L transitions have not been so optimized. The values provided by Krause 
\cite{kraus} have been used. 

For fitting of the spectral peaks, a Gaussian fitting function was used with a width which
expands slightly as the energy level increases, to agree with observed variation.

The identification of a particular element requires a close fit of multiple transitions. 
When the presence of a particular element is to be tested in the presence of overlapping 
fittings, the unknown element’s predicted spectrum is fitted to that portion of the 
measured spectrum which is not already accounted for by the fittings from the other 
elements. The fitting algorithm developed does not allow all sets of elemental peaks to 
be freely variable to fill the available peak intensity envelope. A criterion for 
introducing a new elemental fitting is that the amount of fitted signal will be limited 
by the least-represented line to prevent misattribution of overlapping transitions. Thus, the 
fitting algorithm is ``greedy'' (fits against as much signal as it can) on an 
element-by-element basis, but not on a peak-by-peak basis. Several fitting sequences 
need to be tested to ensure that the solution is consistent. Further discussion of these 
points can be found in a paper discussing Peakaboo \cite{mcintyre}.


\tocchapter{Mapping}

To map the fitted elements, click on \button{map}{Map Fittings} in the toolbar. 
Click on the dropdown arrow to the right of the \button{map}{Map Fittings} button to 
choose either the fitting area (default) or fitting height for mapping. Only those 
fittings which are enabled in your \command{Fitted Elements} list will be mapped. It 
is recommended that you enable all fittings for mapping as each fitting results in 
an individual map which can be disabled later. If your map is quite large, or you 
are using time-intensive filters, then it may be better to map fewer fittings at a 
time as the processing on your desktop or laptop will take considerable time.

Once you have clicked on \button{map}{Map Fittings}, a window will pop up indicating the
progress of mapping the fitted elements.

\screenshot{map-window}{Peakaboo's mapping window}

Once the mapping is complete, a window will appear with a coloured map, scale, and 
list of mapped element lines. For datasets which contain dimension information, map 
widths and heights are automatically set to appropriate values. Otherwise, the user 
must adjust the width and height. All fittings will be represented in the initial map. 
To get the individual element maps, uncheck all other fittings, leaving the element 
fitting you wish to view. If your maps have a pixelated appearance, you may want to 
try interpolating the data and using contouring.

if you wish to scale a fitting to its own highest intensity, select \command{Visible Fittings}
under the \command{Scale Intensities by:} section. If you choose the 
\command{All Fittings}, then the maps will be scaled to the highest intensity from 
the sum of all fittings. A minor element will show very a low intensity 
distribution with this latter option, but appear much more intense with the former.

Element line maps may be saved by clicking on the \button{device-camera}{Save Image} 
button on the toolbar or by selecting \command{Maps \menu\button{device-camera}{Save Image}}
from the menu. Various formats are available for saving your element line maps, 
Pixel Image (\command{PNG}), Vector Image (\command{SVG}) and \command{PDF} format.


\tocsection{Overlays}

The \command{Overlay} option for mapping is chosen from the drop down list under 
\command{Mapped Fittings}. This allows you to sort fittings into three groups: 
\command{Red}, \command{Blue}, \command{Green}. These groups are overlaid, with the
presence of each colour indicating the intensities of their fitting groups.

\screenshot{map-overlay}{Overlay of Iron and Zinc concentration}

The \command{Scale Colours} options control the way that the fitting group's colours are scaled.
With the \command{As a Group} option, the highest intensity from \emphasis{all} 
groups of fittings is used to scale the intensities for \emphasis{every} group. With the 
\command{Separately} option, the highest intensity from \emphasis{within each} 
group of fittings is used to scale the intensities for that group. The \command{Separately}
option produces maps which may be qualitatively interesting, but are quantitatively 
incorrect.

\tocsection{Ratios}

Mapping of element ratios is also available within Peakaboo. Select \command{Ratio} from the
\command{Mapped Fittings} dropdown list. Select the fittings (or fitting sets) you wish to
view in ratio and choose the desired colours from the dropdown lists next to each
fitting.

\screenshot{map-ratio}{Ratio of Iron and Zinc}

\tocsection{Errant Data}

Missing or skewed spectra within a data set can be removed from the map. Hover the mouse
over the bad data point in the mapping window, and the index of that data point will 
appear in the status bar at the bottom of the window. To view the errant data point, 
return to the plotting window, making sure you are in \command{Individual Spectrum} mode, and 
enter the \command{Index \#} of the invalid data in the \command{Scan} field at the bottom 
of the window. 

Click on the \button{choose-cancel}{Exclude} button to flag the scan as bad, and to  exclude it from 
the data set. In all future maps, this data point will be interpolated from neighbouring data.


\tocsection{Examining Subsections}

The distribution of elements will likely vary across your sample and hence your
maps. For samples that have a composition that consists primarily of low atomic 
number (low $Z$) elements, it is possible to make a rough estimate of the composition. 
This estimate is based on the relative intensities of many elements (from $Z=20\ldots 44$) 
measured in a glass matrix. The relative intensities of these elements have been entered 
into Peakaboo as conversion coefficients. Using these coefficients the concentration of 
elements within a low $Z$ matrix can be estimated.

\screenshot{select-subset}{Selecting a subset of a data set}

Left click and drag the mouse in order to draw a rectangle surrounding the region
you wish to select. To view the relative intensity information, click on 
\button{badge-info}{Get Intensities} in the mapping window toolbar.

\screenshot{intensity-estimate}{Rough estimate of elemental intensity in the selected region}

Peakaboo also allows you to view spectra from selected regions within your map. Using the 
same technique as described above to select a subsection of your data, click 
\button{view-subset}{Plot Region} to view the spectra for the selected region. Another 
Peakaboo window will open up to display the spectra from the selected region or subset. 
The spectra from the selected region or `subset' can then be processed with Peakaboo.

\screenshot{plot-subset}{Plotting window showing data for a subset of the original data set}

\tocchapter{Filters}

\tocsection{Using Filters}

Filters can be accessed by clicking on the \command{Filters} tab and then clicking on 
\button{edit-add}{Add Filters} to display the available filters. Filters are grouped into
types. Each type heading can be expanded by clicking on the expander to the left of the filter 
type. Select the filter you wish to use and then click \button{choose-ok}{OK} 
to add the filter. 

Once you have added the filter, if you want to modify or check the 
parameters associated with your filter, click on the \button{misc-preferences}{Settings} 
button next to the filter name to bring up the filter settings window. As you change 
the parameters, you will see the changes reflected in the spectrum being displayed. Once 
you are satisfied with the filter settings, you may close the window.

To view a full listing of built-in filters and their descriptions, see the appendix.

\screenshot{brukner-filter}{The Brukner filter's settings window}

The use of filters will obviously alter the spectrum. If you wish to see the original spectrum
in addition to the filtered data, select \command{View \menu\ Raw Data Outline}.

\screenshot{raw-outline}{Background-subtracted data with the original data drawn in red.}

\tocsection{Filtering Example with Normalizer Filter}

When detector deadtime is very high, the detector cannot keep up with the incoming x-rays.
For example, low energy x-rays such as those coming from \element{Ar} may not get counted
and mapping the distribution of \element{Ar} will reveal areas of next to no intensity. 

The \command{Normalizer Filer} scales each spectrum so that the intensity at a given channel is
always the same. This filter can be used in this case to correct the spectra. By normalizing 
against a channel from within the \element{Ar} peak, we can mitigate the distortion caused by 
detector saturation.

Go to the \command{Filters} tab and add the \command{Normalizer Filter} found under the 
\command{Advanced} filter types. Click on the \button{misc-preferences}{Settings} button 
to bring up the settings window for the filter. Place the cursor on the \element{Ar} peak in the 
spectrum and read the channel number and value (counts) from the information found 
below the spectrum window, and enter them into the appropriate settings fields. Make sure 
you are viewing the \command{Mean Average} spectrum to get the average value (counts) for 
the \element{Ar}. Close the edit window.

\clearpage
\screenshot{normalize-uncorrected-ar}{Uncorrected \element{Ar} with weak signal where \element{Ti} is strongest}
\screenshot{normalize-uncorrected-ti}{Uncorrected \element{Ti}}
\screenshot{normalize-corrected-ar}{Corrected \element{Ar}}
\screenshot{normalize-corrected-ti}{Corrected \element{Ti} with greater maximum intensity}

\tocchapter{Extending Peakaboo}

Peakaboo allows users who are familiar with Java programming to extend its functionality by
creating custom filters and adding support for new file formats. Users who are interested in
developing extensions for Peakaboo will need to make sure they have the Java Development Kit (JDK)
installed in their computer (available from \href{http://java.com}{Java.com}). A good code editor,
such as Eclipse, is also recommended. 

Plugins in Peababoo are classes which extend a specific superclass. The superclass in question 
depends on the kind of plugin being created. These classes should be bundled in \file{.jar} files
and placed in the user's application data directory. For example, on Linux, this path is 
\file{\textasciitilde/.config/Peakaboo}.

In order to know which classes to load without scanning each class in every jar file, a jar 
containing a plugin must include a manifest file, located at \file{META-INF/services/} under
the root of the jar file. The file name must be the qualified name of the superclass that this type
of plugin extends, and the contents of the file must be the qualified names of each class to be 
loaded as a plugin, one per line.

Example plugins for both filters and file formats are available online at 
\href{http://sciencestudioproject.com/}{sciencestudioproject.com}, along with a jar file 
containing Peakaboo for building against, as part of the Peakaboo 4 Software Development Kit.

\tocsection{Filters}

Custom filters should extend the \class{AbstractSimpleFilter} or \class{AbstractBackgroundFilter} 
classes, and should either not define a constructor, or define a \code{public} no-argument constructor
which calls \code{super}, the superclass's constructor.

\tocsubsection{Parameters}

Filters define a series of \class{Parameter}s, each with an associated data type. These parameters are 
displayed in the filter's settings panel, using the associated data types to determine what kind of 
widgets should be used. These parameters should be created and initialized in the filter's \code{initialize} 
method. \class{Parameter}s can be created using the following constructor:

\codeblock{%
\class{Parameter}(\\
\tab \class{String} name,\\
\tab \class{ValueType} type,\\
\tab \class{Object} value\\
)
}

\class{ValueType} is an enumeration containing a listing of the different possible types of data that this 
\class{Parameter} can hold.

\codeblock{%
public static enum \class{ValueType}\\
\{\\
\tab INTEGER, \comment{//Integer}\\
\tab REAL, \comment{//Float or Double}\\
\tab BOOLEAN, \comment{//Boolean}\\
\tab SET\_ELEMENT , \comment{//Object}\\ 
\tab FILTER, \comment{//AbstractFilter}\\
\tab SEPARATOR, \comment{//null}\\
\tab CODE \comment{//String}\\
\}
}

Once constructed, the parameter should be passed to the \code{addParameter} method of the superclass. This 
method returns an integer value which you can use to retrieve the parameter later using 
\code{getParameter(int index)}. The alternative is to simply keep local references to all your \class{Parameter}s.
Here is an example of the \code{initialize} method of the \class{Addition} filter using the integer value to 
track the parameter:

\codeblock{%
@Override\\
public void initialize()\\
\{\\
\tab AMOUNT = addParameter(\\
\tab\tab new \class{Parameter}("Amount to Add", \class{ValueType}.REAL, 1.0)\\
\tab );\\
\}
}

This creates a new \class{Parameter} called ``Amount to Add'' with a value type of \code{REAL} and a 
default value of 1.0, and adds it to this filter's set of parameters.

For choosing from one of several preset values, the \code{SET\_ELEMENT} value type can be used in conjunction with
an alternate constructor for the \class{Parameter} class.

\codeblock{\\
\class{Parameter}(\\
\tab \class{String} name,\\
\tab \class{ValueType} type,\\
\tab \class{Object} value,\\
\tab \class{Object}[] possibleValues\\
)
}

For more complex \class{Parameter} value types, a parameter may have properties associated with it. 
For example, the interface widget used to display \code{CODE} parameters looks for properties named 
\code{``Language''} and \code{``ErrorMessage''}; the former used to control syntax highlighting, and 
latter used during the \class{Parameter} validation process to report compilation/execution errors to 
the user. Currently, the code editor is the only one to make use of properties.


\tocsubsection{Validation}

Whenever the user makes changes to the filter's parameter values in the \button{misc-preferences}{Settings} dialog, 
the new combination of parameter values is validated in a call to the filter's \code{validateParameters}
method. This method returns \code{true} if the new combination of parameter values is acceptable, and \code{false}
otherwise.

If the parameter values do not validate, then the values of the input fields are reset to their old values. 
The only exception to this is the \command{CODE} data type, the editor for which displays an error message
rather than reverting someone's code changes.


\tocsubsection{Processing Spectra}

An active filter's \code{filterApplyTo(\class{Spectrum} data)} method is called repeatedly. This method should 
return a modified \class{Spectrum} containing the results of filtering the input \class{Spectrum}

The \class{Spectrum} class has a \class{List}-like interface, but uses a fixed-size \class{float} array to 
store the actual data for memory and performance reasons.  If you prefer to work with the \class{float[]} data directly, the \code{backingArray}
method provides direct access to it.

\tocsection{File Formats (Data Sources)}

Adding support for new file formats in Peakaboo is accomplished by creating a new data source. When creating a 
new data source, the \class{AbstractDSP} (DSP being short for Data Source Plugin) should be extended. This subclass 
should either not define a constructor, or define a \code{public} no-argument constructor which calles 
\code{super}, the superclass's constructor. If you are attempting to integrate existing code and cannot extend 
\class{AbstractDSP}, the \class{DataSource} interface (which \class{AbstractDSP} implements) can be implemented 
instead, and wrapped in a \class{DelegatingDSP}.


\tocsubsection{Defining File Extensions}

Providing a list of file extensions that your data will typically have allows users to browse for your data format
more easily, and is the first step in filtering out data that your data source does not support. This is done with
the \code{List<String> getFileExtensions()} method that you must implement. Note that the file extensions you 
provide are not binding, and the user may select any file with any file extension.

\tocsubsection{Checking the Format of a File}

Once the user has selected a file or files to open, Peakaboo determines which (if any) data sources are able to 
read the selected data. The data source must implement the \code{boolean canRead(\class{String} filename)} and 
\code{boolean canRead(\class{List}<\class{String}> filenames)} methods, returning \code{true} if the data source believes 
that it is able to read the given file(s). While this method should avoid performing time consuming checks like 
reading in the whole data set just to see if it can, it should still perform enough of a check to differentiate it 
from other file formats which may share the same file extension.

\tocsubsection{Reading the Data}

Once your data source has been selected for reading the data, one of the \code{void read(\class{String} filename)} and 
\code{void read(\class{List}<\class{String}> filenames)} methods will be called. The file names given here will be 
the same as the file names given to the \code{canRead} method earlier, so there is no need to recheck the files.

When reading the data, there are several methods defined by \class{AbstractDSP} which help to provide a smoother 
user experience.

\codeblock{%
boolean isAborted()\\
void haveScanCount(\class{int} scanCount)\\
void newScansRead(\class{int} number)
}

The \code{isAborted} method checks to see if the user has decided to abort the read operation. This method 
should be checked periodically, and if it returns \code{true}, the read operation should be halted. The data source
will not be used after the user has requested the read be aborted, so there is no need to worry about any 
inconsistent state caused by aborting the read.

The \code{haveScanCount} method reports the number of scans in the data set. The \code{newScansRead} method 
reports that a certain number of new scans have been read. These two methods used together allow Peakaboo to
show the progress of the read operation to the user. In some file formats, it is not feasible to determine the 
number of scans in advance of reading them. In this case, simply do not call either of these methods, and the
interface will show a ``busy'' indicator instead.


\tocsubsection{Accessing the Data}

Once the \code{read} method has been called, your data source will need to provide access to the data it 
has been asked to read. the \code{\class{Spectrum} get(\class{int} index)} method should provide the spectral data at the
given index.

The \class{Spectrum} class has \class{List}-like interface, but uses a fixed-size \class{float} array to 
store the actual data. A new \class{Spectrum} can be created using one of 
\code{Spectrum(\class{float[]} data)} which copies the array, or 
\code{Spectrum(\class{float[]} data, \class{boolean} copy)} which optionally does not.

\tocsubsection{Memory Issues with Large Data Sets}

When working with large data sets, memory can sometimes become an issue, especially with older machines or 
when the JVM's heap size is small. The \class{SpectrumList} class provides a \class{List} interface for storing 
\class{Spectrum} objects by writing the data to disk and reading it back in using a very fast serialization library 
called Kryo. This allows your data source to read large data sets without having to worry about memory 
constraints. The only caveat is that modification of data read out of this list will not be present when
the data is read from the list again, unless the list is explicitly updated with the modified data.

The \class{SpectrumList} class implements the \class{List} interface, with the exception that it can act as a sparse
list. With \class{SpectrumList} lists, the following is valid, and will not throw an \class{IndexOutOfBoundsException}

\codeblock{%
\class{List}<\class{Spectrum}> spectrumlist = \class{SpectrumList}.create("My Data Set");\\
spectrumlist.set(4, new Spectrum(10));\\
\class{Spectrum} newspectrum = spectrumlist.get(4);
}

The static method \code{\class{List}<\class{Spectrum}> \class{SpectrumList}.create(String name)} will return 
a new \class{SpectrumList} object. In the event of an \class{IOException} from using the temporary file, another (in-memory)
\class{List} implementation which also acts as a sparse list will be returned instead.


\tocsubsection{Providing Extended Information}

Each data source reads from different kinds of files, which contain different kinds of information about the scan. 
Each data source indicates what kinds of information are available through the 
\code{\class{boolean} hasMetadata()} and \code{\class{boolean} hasScanDimensions()} methods. 

The \code{hasMetadata} method returns \code{true} if it has information such as the name of the scan, the time it was 
collected, the name of the facility it was collected at, and so forth. If this method returns \code{true}, then this 
data source must return values for the methods defined in the \class{DSMetadata} interface. If it returns \code{false}, it
should throw \class{UnsupportedOperationException}s when these methods are called.

The \code{hasScanDimensions} method returns \code{true} if it has information such as the data dimensions, the physical 
dimensions, and so forth. If this method returns \code{true}, then this data source must return values for the methods 
defined in the \class{DSScanDimensions} interface. If it returns \code{false}, it should throw 
\class{UnsupportedOperationException}s when these methods are called.

\cleardoublepage
\appendix


\chapter{Filter Descriptions}
\tocsection{Background Filters}

\tocsubsection{Polynomial}

This filter attempts to fit a series of parabolic (or higher order single-term) curves 
under the data, with a curve centred at each channel, and attempting to make each curve 
as tall as possible while still staying completely under the spectrum. The union of 
these curves is calculated and subtracted from the original data.

\tocsubsection{Brukner}

This filter removes background over several iterations by smoothing the data and 
taking the minimum of the unsmoothed and smoothed data for each channel on each pass.

\tocsubsection{Linear Trim}

This filter examines all pairs of points which are $n$ channels apart (ie (1, 10), 
(2, 11), \ldots\ where $n = 10$). For each pair of points, any signal which exceeds a straight 
line connecting the two points is truncated.

\tocsection{Noise Filters}

\tocsubsection{Savitsky-Golay}

This filter attempts to remove noise by fitting a polynomial to each
point $p_i$ and its surrounding points $p_{i-n} \ldots p_{i+n}$, and then taking the value of the
polynomial at $p_i$. A moving average may be considered a special case of this
filter with a polynomial of order 1.

\tocsubsection{Wavelet Low-Pass}

This filter attempts to reduce high-frequency noise by
performing a Wavelet transformation on the spectrum. This breaks the data down
into sections each representing a different frequency range. The high-frequency
regions are then smoothed, and a reverse transform is applied.

\tocsubsection{Aggressive Wavelet Low-Pass}

This filter attempts to reduce high-frequency noise by performing a wavelet 
transformation on the spectrum. This breaks the data down into sections each 
representing a different frequency range. The high-frequency regions are then 
completely removed, and a reverse transform is applied.

\tocsubsection{Fourier Low-Pass}

This filter transforms the spectral data with a Fourier Transformation into a 
frequency domain. Data from a high frequency range (noise) is filtered out, 
while lower frequencies (peaks, background) are passed through.

\tocsubsection{Moving Average}

This filter refines the values of each point in a scan by sampling it and the 
$n$ points to either side of it, and replacing it with an average of the sampled points.

\tocsubsection{Spring Smoothing}

This filter operates on the assumption that weak signal should be smoothed more 
than strong signal. It treats each pair of adjacent points as if they were
connected by a spring. With each iteration, a tension force draws neighbouring
points closer together.

The \command{Force Multiplier} controls how strongly a 
pair of points are pulled together, and the \command{Force Falloff Rate} controls 
how aggressively stronger signal is anchored in place, unmoved by tension forces. 
This prevents stronger intensity such as peak shapes from being distorted by the smoothing 
algorithm.

\tocsection{Mathematical Filters}

\tocsubsection{Add}

This filter adds a constant value to all points on a spectrum.

\tocsubsection{Subtract}

This filter subtracts a constant value to all points on a spectrum.

\tocsubsection{Multiply}

This filter multiplies all points on a spectrum by a constant value.

\tocsubsection{Derivative}

This filter transforms the data such that each channel represents the
difference between itself and the channel before it.

\tocsubsection{Integral}

This filter transforms the data such that each channel represents the sum of
itself and all channels prior to it.


\tocsection{Advanced Filters}

\tocsubsection{Signal $\rightarrow$ Wavelet}

This filter converts spectrum data into a wavelet representation. This is intended
to be used in conjunction with other filters (especially the \command{Filter Partial Spectrum}
filter) to perform custom wavelet operations.

\tocsubsection{Wavelet $\rightarrow$ Signal}

This filter converts a wavelet representation of data back into
spectrum data. This is intended to be used in conjunction with other filters
(especially the \command{Filter Partial Spectrum} filter) to perform custom 
wavelet operations.

\tocsubsection{Normalizer}

This scales each spectrum so that the intensity at a specified channel is
always the same across all spectra.

\tocsubsection{Filter Partial Spectrum}

This filter allows the application of another filter to a portion of a spectrum.


\tocsection{Programming Filters}

\tocsubsection{Java Code}

This filter allows you to create your own filter in the Java programming language.

\tocsubsection{Python Code}

This filter allows you to create your own filter in the Python programming language using Jython.




\chapter{File Format Descriptions}

\tocsection{CDFML Files}

CDFML XRF files are a format defined by the Canadian Light Source using NASA's CDF/CDFML container.

\tocsection{Peakaboo Plain Text Files}

Peakaboo Plain Text format is a simple XRF format comprised of rows of space-separated numbers. Each row 
represents a single spectrum.

\tocsection{CLS Data Acquisition Format}

Data format used by the Canadian Light Source for XRF collection on the VESPERS beamline

\tocsection{MCA Files}

Experimental support for the MCA XRF data format.

\bibliographystyle{plain}
\bibliography{manual}

\end{document}

